\documentclass[11pt,a4paper]{article}

\usepackage[T1]{fontenc}
\usepackage[utf8]{inputenc}
\usepackage{lmodern}
\usepackage[a4paper,margin=26mm]{geometry}
\usepackage{amsmath,amssymb,amsthm,mathtools}
\usepackage{booktabs,array}
\usepackage{microtype}
\usepackage{xcolor}
\usepackage{enumitem}
\usepackage{url}
\usepackage[colorlinks=true,linkcolor=blue!55!black,citecolor=blue!55!black,
            urlcolor=blue!55!black]{hyperref}

\newtheorem{theorem}{Theorem}
\newtheorem{lemma}[theorem]{Lemma}
\newtheorem{proposition}[theorem]{Proposition}
\newtheorem{remark}[theorem]{Remark}
\newcommand{\F}{\mathbb F}
\newcommand{\bits}{\{0,1\}}
\newcommand{\sample}{\mathrel{\leftarrow\$}}
\newcommand{\CF}{\mathsf{CF}}
\newcommand{\Corank}{\mathsf{Corank1Cal}}
\newcommand{\Parse}{\mathsf{ParseHash}}
\newcommand{\Sign}{\mathsf{Sign}}
\newcommand{\Verify}{\mathsf{Verify}}
\newcommand{\Prb}{\mathop{\rm Pr}}
\newcommand{\qfld}{q_{\mathrm{fld}}}

\title{\bfseries Ordinary Short-Seed Collisions in the Unsalted\\
TRINE Specification:\\
Equivalent Signing-Witness Recovery and Fresh-Message Forgery}
\author{Martin Feussner\\
\small Selmer Center, University of Bergen\\
\small \texttt{martin.feussner@uib.no}}
\date{3 October 2026}

\begin{document}
\maketitle

\begin{center}
\fbox{\parbox{0.92\linewidth}{\small\textbf{Provenance and review status.}
The TRINE-specific attack instantiation, derivation, implementation,
experiments, resource analysis, and initial manuscript were produced by OpenAI
Codex using Daybreak Blue at Ultra reasoning effort during an AI-run audit of
NGCC signature candidates.  A separately instructed hostile-review agent
reconstructed the attack, corrected the resource accounting, rebuilt the
full-dimension reduced-seed experiments, and tried to disprove the result.  At
the date above, no human cryptanalyst has independently verified the argument,
implementation, measurements, or novelty assessment.  This report is a
preliminary disclosure for human review and reproduction.}}
\end{center}
\vspace{0.7em}

\begin{abstract}
The submitted TRINE specification compresses every base opening in a signature
to a $\lambda$-bit seed.  Algorithms~6--8 and the signature-size formula contain
no per-signature salt, and the point sampled from that seed is fixed before the
Fiat--Shamir hash sees the message.  Under the direct, message-independent
construction described in the specification, ordinary independently randomized
signatures eventually reuse a round seed at the same round index.  If the two
Fiat--Shamir labels differ, the public canonical forms are equal.  The public
canonicalization certificates then give a complete equivalence between the two
corresponding trilinear forms.

For level-I Balanced, one base/nonbase collision recovers an equivalent signing
witness.  For level-I ShortSig, collision edges between the five form labels
compose; any connected label graph recovers all four maps from the base form.
The recovered first-coordinate maps suffice to answer every future challenge,
so the attacker signs arbitrary fresh messages without the original secret key.
Full-dimension experiments reduced only the round-seed entropy to make natural
birthday collisions executable.  They used fresh independent DRBG output,
recovered and checked the complete public tensor equivalences, erased the
original secret material, and produced fresh-message signatures accepted by a
specification-conformant verifier.  Wrong-message, response-tamper,
digest-tamper, and corrupted-map controls rejected.

Most directly, set the signing-query count to
$Q=2^{64}-1<2^{64}$.  In the exact fixed-index ideal model, the least public
tables giving success at least one half contain $2^{64.879318}$ and
$2^{67.639106}$ distinct seeds, with one 256-bit canonical fingerprint per
seed, for Balanced-I and ShortSig-I.  The corresponding returned traffic is
$2^{78.609179}$ and $2^{77.661778}$ bits.  A raw 48-byte seed/fingerprint
table occupies $2^{70.464281}$ and $2^{73.224068}$ bytes; conservative 80-byte
records occupy $2^{71.201246}$ and $2^{73.961034}$ bytes.  Measured
full-canonical-form accounting gives $2^{94.246627}$ and $2^{96.990188}$
attacker cycles, while the honest signing service performs $2^{97.268224}$
and $2^{96.182868}$ cycles.  Thus the attack algorithms have full-parameter
cost estimates below the claimed 128-bit level and need fewer than $2^{64}$
signing queries, but they are not physical $2^{80}$ attacks.  Only the
reduced-seed end-to-end forgeries were executed; no full-entropy forgery was
run.

The untouched submitted C source serializes a fresh $2\lambda$-bit salt and
absorbs it in every round expansion.  That source-only field is absent from the
specification algorithms and size formula, makes level-I signatures 32 bytes
longer, and blocks this cross-signature accumulation.  We therefore claim an
attack only on the direct unsalted, message-independent construction described
in the specification, not on the untouched salted source.  The generic
short-seed collision attack and the salt-plus-round repair are due to Beullens
and Sipasseuth.  Our contribution is
limited to the TRINE instantiation: the specification/source discrepancy,
canonical-certificate extraction, one-edge Balanced completion, connected-graph
ShortSig completion, exact level-I resource analysis, and end-to-end validation.
\end{abstract}

\section{Introduction}

TRINE is a Fiat--Shamir signature based on trilinear-form equivalence
(TFE)~\cite{trine-spec}.  Its identification protocol commits to canonical
forms computed from effective corank-one points.  A base challenge opens the
point itself; a nonbase challenge opens the point transported through a secret
equivalence.  The canonical form is both an invariant and an extractor: two
different openings of one commitment reveal an equivalence between the public
forms.

The specification reduces signature size by replacing each base opening with a
$\lambda$-bit seed.  This creates a finite commitment-randomness domain.  With
many ordinary signatures, a seed repeats.  When the Fiat--Shamir labels at that
round differ, the special-soundness event needed to recover a public-form
equivalence occurs naturally.  Balanced needs one such event.  ShortSig has
four nonbase forms, but its recovery problem is only connectivity of a
five-vertex graph: edges from different round indices compose because the
secret equivalences are profile-wide.

Short-seed collision attacks of this kind are known.  Beullens described a
public-table attack on ALTEQ and the need to include a salt and round identifier
in the expansion domain~\cite{beullens2024}.  Sipasseuth's DVA analysis
explicitly discusses both attacker-generated seed tables and internal
collisions, together with the same repair~\cite[Sec.~3.2]{sipasseuth2024}.
This report does not claim those generic ideas.  It establishes that the
unsalted algorithm and size formula in the TRINE specification instantiate the
attack, derives the exact extraction and graph completion for both level-I
profiles, and accounts for the unusually expensive TRINE seed decoder.

Our contributions are as follows.
\begin{enumerate}[leftmargin=*,itemsep=2pt]
  \item We give the exact same-index collision and canonical-certificate
        extraction equations for the construction described in the
        specification.
  \item We prove that one collision edge is a complete Balanced signing
        witness and that any connected collision graph is a complete ShortSig
        signing witness.
  \item We give an exact fixed-index ideal-model attack point at
        $Q=2^{64}-1$ and, separately, derive approximate all-round
        sparse-Poisson sensitivity estimates and broader time--memory--data
        tradeoffs, while accounting separately for signing
        service, seed decoding, canonicalization, returned traffic, retained
        storage, and peak RAM.
  \item We implement the full extraction-to-forgery chain at the exact level-I
        dimensions, reducing only seed entropy.  A second implementation based
        on the submitted source uses ordinary fresh DRBG output and exact
        specified signature sizes.
  \item We isolate the specification/source discrepancy.  The untouched source
        already contains the fresh $2\lambda$-bit salt needed to defeat
        accumulation; it is outside the attack claim and is not vulnerable to
        this mechanism.
\end{enumerate}

The attack is unconditional with respect to signer state: it uses neither
rollback nor repeated-state control, faults, related keys, rare keys, nor a
chosen-randomness interface.  Its scope does depend on the specification's seed
decoder being deterministic and message-independent.  That is the direct
reading of the written sampling order and matches the submitted decoder after
removing only its source-only salt.  A hypothetical decoder keyed by the
message would define a different construction and defeat cross-message
accumulation.

\section{The relevant TRINE construction}

\subsection{Public forms and the group action}

Let $\phi:\F_{\qfld}^n\times\F_{\qfld}^n\times\F_{\qfld}^n
\rightarrow\F_{\qfld}$ be a trilinear form.  For
$T=(A,B,C)\in\mathrm{GL}(n,\qfld)^3$, write
\begin{equation}
 \phi^T(x,y,z)=\phi(Ax,By,Cz).                         \label{eq:action}
\end{equation}
The action is on the right:
\begin{equation}
                    (\phi^T)^S=\phi^{TS}.              \label{eq:rightaction}
\end{equation}

TRINE publishes $X+1$ equivalent forms.  We use label $X$ for the base form:
\begin{equation}
  \phi_j=\phi_X^{E_j},\qquad
  E_j=(A_j,B_j,C_j),\quad 0\le j<X,qquad E_X=(I,I,I).
                                                               \label{eq:publicforms}
\end{equation}
The original secret key contains the nonbase equivalences $E_j$.  Any other
triples satisfying~\eqref{eq:publicforms} are equivalent signing witnesses;
recovering the original key representation is unnecessary.

For an effective corank-one point $a$ of a form $\phi$, TRINE computes the
canonical form $\CF(\phi,a)$.  The public canonicalization procedure also
constructs invertible bases.  We denote its certificate by
$T_{\phi,a}=(P,Q,R)$, so that
\begin{equation}
                       \phi^{T_{\phi,a}}=\CF(\phi,a).  \label{eq:certificate}
\end{equation}

\subsection{Signing and verification equations}

Algorithm~6 first samples all round points and commitments, and only then
hashes the message:
\begin{align}
 a_i&\sample\Corank(\phi_X),\nonumber\\
 \psi_i&=\CF(\phi_X,a_i), &&0\le i<r,                  \label{eq:commitments}\\
 \mathit{cha}&=H(M\mathbin\Vert\psi_0\mathbin\Vert\cdots
                    \mathbin\Vert\psi_{r-1}),\nonumber\\
 (b_0,\ldots,b_{r-1})&=\Parse(\mathit{cha}),\nonumber\\
 d_i&=A_{b_i}^{-1}a_i.                                 \label{eq:response}
\end{align}
The parser chooses exactly $K$ nonbase positions.  Conditional on a nonbase
position, its label is uniform in $\{0,\ldots,X-1\}$; all remaining positions
have label $X$.

Verification reconstructs $\CF(\phi_{b_i},d_i)$, rehashes those forms with the
message, reparses the challenge, and accepts only if the challenge and every
opening are consistent.  Correctness follows from
\begin{equation}
 \CF(\phi_j,A_j^{-1}a)=\CF(\phi_X,a).                  \label{eq:invariance}
\end{equation}

Section~3.4 permits a base point $a_i$ to be sent as a seed.  Section~3.5
assigns exactly $\lambda$ bits to each of the $r-K$ base seeds and gives
\begin{equation}
 |\sigma|=(r-K+2)\lambda+K n\lceil\log_2 \qfld\rceil.   \label{eq:size}
\end{equation}
There is no salt in Algorithms~6--8 or~\eqref{eq:size}.  A conservative direct
realization preserves round separation through public maps
\begin{equation}
 D_i:\bits^\lambda\longrightarrow
       \{\text{effective corank-one points of }\phi_X\},
 \qquad \psi_i(s)=\CF(\phi_X,D_i(s)).                  \label{eq:decoder}
\end{equation}
The attack compares seeds only at equal indices $i$; it never assumes a common
decoder across round indices.  Injectivity is unnecessary: equal inputs always
give equal commitments, while additional decoder collisions only help.

\begin{center}
\fbox{\parbox{0.92\linewidth}{\small\textbf{Normative decoder evidence.}
The direct message-independent reading of $D_i$ rests on cumulative evidence
in the submission.  Algorithm~6 samples $a_i\sample\Corank(\phi_X)$ and
computes every $\psi_i$ before hashing the message, and the point sampler has
no message argument.  Section~3.4 permits a base opening to be compressed to a
seed; Section~3.5 assigns exactly $\lambda$ bits to each such seed.
Algorithms~6--8 and the specification's signature-size formula serialize no
salt.  In contrast, the submitted C source adds a fresh $2\lambda$-bit salt,
absorbs it in every round expansion, and consequently produces level-I
signatures 32 bytes longer.}}
\end{center}

\subsection{Level-I parameters}

Table~\ref{tab:params} gives the two targeted profiles.  The signature sizes
specified by the construction follow directly from~\eqref{eq:size}.  The
submitted source sizes include one extra $2\lambda$-bit field.

\begin{table}[ht]
\centering
\caption{TRINE level-I parameters and signature sizes.}
\label{tab:params}
\begin{tabular}{lrrrrrr}
\toprule
profile & $n$ & $\qfld$ & $r$ & $K$ & $X$ & specification/source bytes\\
\midrule
Balanced-I & 22 & 4093 & 138 & 52 & 1 & 3,124 / 3,156\\
ShortSig-I & 22 & 4093 & 61 & 36 & 4 & 1,620 / 1,652\\
\bottomrule
\end{tabular}
\end{table}

An exhaustive search of the complete 23-page specification, its EUF--CMA and
parameter-selection sections, the basic-information form, both IP documents,
the implementation README and interface notes, all reference and optimized
source trees, and the KAT material found no maximum signatures per key and no
mandatory key-rotation rule.  In particular, KeyGen has no lifetime input,
Section~5.2 states EUF--CMA security without a query cap, and Section~6.3
assigns the 128-bit level from $(r,K,X)$ without a lifetime condition.  This
establishes only the absence of a specified cap; it does not attribute an
affirmative unlimited-lifetime policy to the submitters.  The
$Q=2^{64}-1$ point below is therefore not excluded by a submitted per-key
lifetime rule.  A deployment enforcing a lower cap cannot reach that
particular point.  Its
other resources remain vastly beyond physical execution on the audit machine.

\section{Collision extraction}

\begin{lemma}[Same-index seed collision]\label{lem:seedcollision}
Suppose two signatures use the same underlying round seed $s$ at one round
index $i$, but their challenge labels there are $c$ and $d$.  Under
\eqref{eq:decoder}, their verifier-reconstructed canonical forms are equal:
\[
             \CF(\phi_c,u_c)=\psi_i(s)=\CF(\phi_d,u_d).
\]
This holds for ordinary independent signing randomness; it does not require
the signer state to repeat.
\end{lemma}
\begin{proof}
Both signing calls evaluate the same deterministic map $D_i$ on $s$, giving
the same base point and commitment.  Correctness~\eqref{eq:invariance} carries
that commitment through either challenge response.  Repetition arises from
the finite $\lambda$-bit seed space under independent sampling.
\end{proof}

\begin{lemma}[Public certificate extraction]\label{lem:extract}
Let $c\ne d$ and suppose public openings $u_c,u_d$ satisfy
$\CF(\phi_c,u_c)=\CF(\phi_d,u_d)$.  Let $T_c,T_d$ be the public certificates
from~\eqref{eq:certificate}.  Then
\begin{equation}
                         \phi_c^{T_cT_d^{-1}}=\phi_d.   \label{eq:edge}
\end{equation}
\end{lemma}
\begin{proof}
The equal canonical form gives $\phi_c^{T_c}=\phi_d^{T_d}$.  Applying
$T_d^{-1}$ on the right and using~\eqref{eq:rightaction} proves the claim.
\end{proof}

The attacker detects a candidate edge by comparing canonical-form encodings,
not by observing hidden seed values.  Before accepting an edge, the reproducer
byte-compares the complete canonical tensors and checks~\eqref{eq:edge} against
all public tensor coefficients.  Stabilizers or nonunique certificates do not
matter: any extracted triple that passes this public equation is a valid
equivalence.

\subsection{Balanced: one edge}

Balanced has labels $\{0,1\}$ with base label $X=1$.  A collision with
$c=1,d=0$ yields
\[
 E'_0=T_1T_0^{-1},\qquad \phi_1^{E'_0}=\phi_0.
\]
Let $A'_0$ be its first component.  On every future nonbase challenge the
attacker returns $(A'_0)^{-1}a_i$; on a base challenge it serializes the seed.

\begin{theorem}[Balanced completion]\label{thm:balanced}
One observable same-index collision under the two different Balanced labels
gives an equivalent signing witness and permits valid signatures on arbitrary
fresh messages.
\end{theorem}
\begin{proof}
Lemma~\ref{lem:extract} produces a publicly checked triple mapping the base
form to the only nonbase form.  Equation~\eqref{eq:invariance} shows that its
inverse first component answers every nonbase response.  All other signing
steps use public algorithms and fresh attacker-chosen round seeds.  The
Fiat--Shamir hash therefore binds a correctly opened commitment vector to the
new message.
\end{proof}

\subsection{ShortSig: connected-graph completion}

ShortSig has labels $V=\{0,1,2,3,4\}$ with base label $4$.  Each collision
between distinct labels $c,d$ yields the directed equivalence
$E_{c\rightarrow d}=T_cT_d^{-1}$ and its inverse.  For a path
\[
 4=v_0,v_1,\ldots,v_m=j,
\]
right-action composition gives
\begin{equation}
 E'_{4\rightarrow j}=E_{v_0\rightarrow v_1}
                      E_{v_1\rightarrow v_2}\cdots
                      E_{v_{m-1}\rightarrow v_m},
 \qquad \phi_4^{E'_{4\rightarrow j}}=\phi_j.           \label{eq:path}
\end{equation}

\begin{theorem}[ShortSig completion]\label{thm:shortsig}
If the graph of observed cross-label seed collisions is connected, the
attacker recovers an equivalent signing witness for all four ShortSig nonbase
forms and can sign arbitrary fresh messages.
\end{theorem}
\begin{proof}
Connectivity supplies a path from base vertex $4$ to every nonbase vertex.
Equation~\eqref{eq:path} produces and publicly verifies every required
equivalence.  Inverting the first component of each triple supplies the
response map in~\eqref{eq:response}.  The rest follows as in
Theorem~\ref{thm:balanced}.
\end{proof}

Edges may come from different round indices.  Round separation changes which
seed values collide, but it does not change the profile-wide public-form maps
recovered from those edges.

\section{Collision acquisition and success probabilities}

\subsection{All-round internal collisions}

Process ordinary signatures once, reconstruct every round canonical form, and
compare records only within the same index.  Let
\begin{equation}
 p_j=\frac{K}{rX}\quad(j<X),\qquad p_X=\frac{r-K}{r}.   \label{eq:labelprob}
\end{equation}
For $Q$ signatures, the sparse-Poisson edge intensity between labels $u,v$ is
\begin{equation}
             \mu_{uv}=\frac{Q(Q-1)r p_up_v}{2^{128}}.   \label{eq:edgepoisson}
\end{equation}

For Balanced, success means at least one edge, so
\begin{equation}
 \Prb[\mathrm{success}]\approx 1-exp\!\left(
  -\frac{Q(Q-1)r p_0p_1}{2^{128}}\right).              \label{eq:balprob}
\end{equation}
Its 50-percent point is $\log_2Q=61.226527$.

For ShortSig, assign independent edge probabilities
$1-e^{-\mu_{uv}}$ in the sparse approximation and sum the probabilities of
the connected graphs among the $2^{\binom52}=2^{10}$ edge sets.  The resulting
median is $\log_2Q=62.990592$.  Nonbase/nonbase edges materially improve on a
base-star-only collection.

These all-round probabilities are an ideal sparse-Poisson model, not exact
probabilities for conditioned serialized signatures.  The labels within one
signature have fixed weight and are correlated.  An exact fixed-weight toy
simulation with 16-bit seeds obtained 968 successes in 2,000 Balanced trials
at 37 signatures, versus model probability 0.49183, and 973 successes in 2,000
ShortSig trials at 127 signatures, versus 0.49811.  Digest collisions and
random-oracle conditioning are a separate negligible layer in this model.

\subsection{Fixed-index public table}

Fix one round $i$.  Decode and canonicalize $P$ distinct public seeds under the
base form, then inspect the opening at index $i$ in each of $Q$ ordinary
signatures.  A nonbase opening whose canonical form appears in the table gives
an edge.

Let $N=2^{128}$ and $p=K/r$.  For Balanced, the exact uniform-seed,
ideal-challenge probability is
\begin{equation}
                    1-\left(1-\frac{pP}{N}\right)^Q.   \label{eq:fixedbal}
\end{equation}
For ShortSig, requiring a hit for all four nonbase labels gives by
inclusion--exclusion
\begin{equation}
 \sum_{s=0}^{4}(-1)^s\binom4s
       \left(1-\frac{s pP}{4N}\right)^Q.               \label{eq:fixedshort}
\end{equation}
These are exact for uniform seeds and independent ideal labels at the selected
index; they remain analytic models for the concrete Fiat--Shamir transcript.

At sparse 50-percent points,
\begin{equation}
 QP=\frac{zXN}{p},\qquad
 z=\ln2\ \text{for Balanced},\qquad
 z=-\ln(1-2^{-1/4})\ \text{for ShortSig}.              \label{eq:pareto}
\end{equation}
This exposes the full time--memory--data tradeoff.  If decoding one public seed
costs $c_D$ and processing one nonbase target costs $c_F$, measured attacker
work is
\begin{equation}
                  W(P)=c_DP+pQc_F.                     \label{eq:weighted}
\end{equation}

\begin{proposition}[A fixed-index point below $2^{64}$ signing queries]
\label{prop:q64}
Set
\[
 Q=2^{64}-1=18,446,744,073,709,551,615.
\]
For Balanced-I and ShortSig-I, respectively, set
\begin{align*}
 P_B&=33,932,896,019,960,893,868,\\
 P_S&=229,826,321,259,624,867,428.
\end{align*}
Substitution in the exact ideal-model
formulae~\eqref{eq:fixedbal} and~\eqref{eq:fixedshort} gives, respectively,
\begin{align*}
 \Prb_B[\mathrm{success}]&=
  0.5000000000000000000044518893\ldots,\\
 \Prb_S[\mathrm{success}]&=
  0.5000000000000000000011550258\ldots.
\end{align*}
Thus both specified witness-recovery routes have approximately 50-percent
success with strictly fewer than $2^{64}$ signing queries.
\end{proposition}
\begin{proof}
Use $p=52/138$ in~\eqref{eq:fixedbal}, and $p=36/61$ in
\eqref{eq:fixedshort}.  The companion script evaluates those formulae with
120 decimal digits and integer binary search.  It additionally obtains
$0.4999999999999999999942383907\ldots$ at $P_B-1$ and
$0.4999999999999999999981283883\ldots$ at $P_S-1$.  Hence the displayed
integers are the least table sizes reaching one half under the stated exact
fixed-index ideal model.  No sparse-Poisson approximation enters this result.
For ShortSig the four hits make a base-centered star, which is connected, so
Theorem~\ref{thm:shortsig} applies.
\end{proof}

This point matters independently of its large physical work: it does not
require $2^{64}$ signing queries, and the submission specifies no lower
per-key lifetime cap.  It uses ordinary independently randomized signatures.

\section{Concrete resources}

Every exponent in this section is base two.  The figures use the level-I
signature sizes given by the specification.  ``Signer cycles'' charge all signing-oracle service
separately using the cycle counts in the submitted performance table.
``Attacker cycles'' use native measurements on an AMD EPYC host at 1.996249
GHz.  Parallelism changes wall time, not aggregate work.

\subsection{Measured primitive costs}

The decoder repeatedly samples field points until $\Corank$ and canonicalization
succeed.  A success occurs only about once per 4,093 candidates.  The hostile
benchmark measured the complete seed-to-canonical-form path:
\begin{center}
\small
\begin{tabular}{lrr}
\toprule
operation & Balanced-I & ShortSig-I\\
\midrule
seed $\rightarrow\Corank$/retry $\rightarrow\CF$, seconds &
  0.345799062 & 0.342778750\\
seed $\rightarrow\Corank$/retry $\rightarrow\CF$, cycles &
  $2^{29.36265}$ & $2^{29.34999}$\\
full $\CF$ of an exposed point & \multicolumn{2}{c}{$0.005459163$ s
  ($2^{23.37754}$ cycles)}\\
12-value partial invariant & \multicolumn{2}{c}{$0.000599711$ s
  ($2^{20.19120}$ cycles)}\\
\bottomrule
\end{tabular}
\end{center}
The seed measurement already includes the full canonical form.  Omitting this
retry path understates the attack by many bits.

\subsection{The \texorpdfstring{$Q=2^{64}-1$}{Q below 2 to the 64} security point}

Table~\ref{tab:q64} accounts for Proposition~\ref{prop:q64}.  Each public-table
entry contains one distinct 16-byte seed and one 32-byte canonical-form
fingerprint.  The raw table therefore uses 48 bytes per entry.  The second
storage row conservatively allows 80 bytes for sorting and retrieval metadata,
as in the all-round accounting.  ``Query-side hashes'' is the expected number
$pQ$ of nonbase openings at the selected index that must be canonicalized and
fingerprinted.  On a fingerprint match, the attacker regenerates and
byte-compares both complete 15,972-byte canonical tensors and accepts an edge
only after checking the public tensor-equivalence equation coefficient by
coefficient.

\begin{table}[ht]
\centering
\small
\setlength{\tabcolsep}{5pt}
\caption{Exact fixed-index ideal-model points with fewer than $2^{64}$
signing queries.  Resource entries other than $Q$ and success are base-two
exponents.}
\label{tab:q64}
\begin{tabular}{lrr}
\toprule
resource & Balanced-I & ShortSig-I\\
\midrule
signing queries $Q$ & $2^{64}-1$ & $2^{64}-1$\\
public table seeds & $2^{64.879318}$ & $2^{67.639106}$\\
public table hashes & $2^{64.879318}$ & $2^{67.639106}$\\
expected query-side hashes & $2^{62.591915}$ & $2^{63.239188}$\\
returned signature bits & $2^{78.609179}$ & $2^{77.661778}$\\
raw table bytes (48/entry) & $2^{70.464281}$ & $2^{73.224068}$\\
operational table bytes (80/entry) & $2^{71.201246}$ & $2^{73.961034}$\\
measured attacker cycles & $2^{94.246627}$ & $2^{96.990188}$\\
honest signer cycles & $2^{97.268224}$ & $2^{96.182868}$\\
exact ideal success & $0.5+4.45\mathbin{\cdot}10^{-21}$ &
                      $0.5+1.16\mathbin{\cdot}10^{-21}$\\
\bottomrule
\end{tabular}
\end{table}

The full-canonical attacker count is
\[
 P(c_D\cdot1.996249\cdot10^9)
 +pQ(c_F\cdot1.996249\cdot10^9)
\]
cycles, with $c_D$ and $c_F$ from the measured table above.  The signer count
is $Q$ times the submission's reported signing cycles.  These are aggregate
work figures; parallel execution changes wall time but not the displayed
cycle totals.  The exact probability belongs to the uniform-seed,
ideal-challenge model at one fixed index.  It is not asserted as the exact
probability of a conditioned serialized Fiat--Shamir transcript.

\subsection{All-round route}

Each signature yields $r$ externally sortable records.  The attacker decodes
every base seed and canonicalizes each exposed nonbase point.  Table
\ref{tab:allround} reports the median point.

\begin{table}[ht]
\centering
\caption{All-round internal-collision resources.}
\label{tab:allround}
\begin{tabular}{lrr}
\toprule
resource & Balanced-I & ShortSig-I\\
\midrule
signing queries & $2^{61.226527}$ & $2^{62.990592}$\\
returned signature bits & $2^{75.835706}$ & $2^{76.652370}$\\
round records / canonical results & $2^{68.335051}$ & $2^{68.921330}$\\
retained bytes at 80 bytes/record & $2^{74.656979}$ & $2^{75.243258}$\\
measured attacker cycles & $2^{97.029148}$ & $2^{97.017155}$\\
honest signer cycles & $2^{94.494751}$ & $2^{95.173460}$\\
\bottomrule
\end{tabular}
\end{table}

The $80$-byte record contains a canonical digest, label, round index, opening or
seed, and retrieval metadata.  External sorting permits an operator-selected
RAM chunk; peak RAM need not approach the retained-disk exponent.  On a digest
match, the attacker regenerates and byte-compares both 15,972-byte canonical
forms before extraction.  A record-hash collision can add work but cannot
produce a false witness.

\subsection{Fixed-index Pareto points}

Table~\ref{tab:fixed} separates three useful points.  The first minimizes the
number of fingerprint evaluations and keeps raw traffic just below $2^{80}$
bits.  The second minimizes measured aggregate attacker cycles with full
canonical forms and is the rigorous CPU baseline.  The third holds signing
queries at the all-round medians, minimizing traffic but requiring a much
larger public table.

\begin{table}[ht]
\centering
\caption{Fixed-index full-canonical-form Pareto points.  Columns headed $B$
and $S$ are Balanced-I and ShortSig-I.}
\label{tab:fixed}
\begin{tabular}{lrrrrrr}
\toprule
& \multicolumn{2}{c}{minimum calls} & \multicolumn{2}{c}{minimum attacker CPU}
& \multicolumn{2}{c}{all-round $Q$}\\
resource & $B$ & $S$ & $B$ & $S$ & $B$ & $S$\\
\midrule
$\log_2 Q$ queries &65.1437&66.2000&68.1363&69.1862&61.2265&62.9906\\
$\log_2 P$ table seeds &63.7356&65.4391&60.7430&62.4529&67.6528&68.6485\\
returned bits &79.7529&79.8618&82.7455&82.8480&75.8357&76.6524\\
table bytes &70.0575&71.7610&67.0649&68.7748&73.9747&74.9704\\
attacker cycles &93.1209&94.8119&91.1057&92.8029&97.0155&97.9988\\
signer cycles &98.4119&98.3829&101.4045&101.3691&94.4948&95.1735\\
\bottomrule
\end{tabular}
\end{table}

At the CPU minimum, the public seed table is much smaller but the oracle
returns more data and performs roughly $2^{101.4}$ cycles of signing work.  At
the data-oriented point, attacker work rises to $2^{93.12}$ and $2^{94.81}$
cycles.  No row is a simultaneous $2^{80}$ time/data/memory attack.

\subsection{Partial invariant acceleration}

The canonicalizer's alternating kernel chain produces vectors that transform
up to independent nonzero scales.  Let
$T_{ijk}=\phi(u_i,v_j,w_k)$ and $d=T_{010}$.  When the denominators are
nonzero, the ratios
\begin{align}
 \frac{dT_{ij0}}{T_{i10}T_{0j0}}&& (i,j\in\{2,3\}),\nonumber\\
 \frac{d^2T_{ijk}}{T_{i10}T_{0j0}T_{01k}}&&
   (i,j,k\in\{2,3\})                                  \label{eq:ratios}
\end{align}
cancel every chain-vector scale.  The resulting 12 elements of $\F_{4093}$
have a nominal 143.987 bits of output.

An independent full-dimension primitive test obtained equality for 64/64 planted
equivalences and 0/64 independent controls, with a 9.10-fold speedup over full
canonicalization.  Equivariance is exact, but the tuple's generic
discrimination entropy on TRINE's structured distribution is not proved.
Every hit must therefore receive full-$\CF$ confirmation.  Treating the tuple
as a high-entropy prefilter gives $2^{89.51}$ and $2^{91.21}$ measured cycles
at its CPU-optimized points; these are heuristic sensitivity figures.  The
$2^{91.11}$ and $2^{92.80}$ full-$\CF$ rows remain the rigorous fixed-index
baselines.

\section{End-to-end validation}

\paragraph{Full-entropy specification conformance.}
A clean copy of the submitted reference source received only the minimal
salt-removal and signature-layout patch described below.  The round seed
remained 16 bytes (128 bits); no collision instrumentation or entropy
reduction was compiled.  One ordinary Balanced-I signature verified and
opened at exactly 3,124 bytes, and one ordinary ShortSig-I signature verified
and opened at exactly 1,620 bytes.  Changed-message controls rejected.  This
is a full-entropy conformance test of the direct unsalted transcript defined by
the specification, not a collision, extraction, or forgery experiment.  It
establishes that reducing seed entropy is needed only to make the birthday
events experimentally reachable.

\paragraph{Reduced-seed attack validation.}
Two implementations tested the complete extraction-to-forgery chain while
retaining the exact level-I dimensions $(n,\qfld,r,K,X)$ and reducing only
seed entropy.

\paragraph{Portable tests.}
With eight-bit independent fresh seeds, Balanced needed four ordinary
signatures and ShortSig needed nine in the frozen runs.  Both tests detected
collisions through public canonical equality, recovered the required tensor
maps, erased the planted signing maps, and forged a fresh message using only
the recovered witness.  Wrong-message and corrupted-map controls rejected.

\paragraph{Submitted-source hostile harness.}
A separately written harness aligns the submitted level-I core with the
transcript defined by the specification, retains round-index separation, and
reduces the seed domain to six
bits.  For instrumentation convenience, the former salt input is fixed to the
public all-zero string and omitted from serialization; a fixed public value
adds no per-signature decoder entropy.  The bundle separately includes the
minimal conformance patch that removes the source-only salt from the decoder
as well.  The submitted DRBG is explicitly initialized and never reset or
rolled back.  Across two independently generated keys per profile, Balanced
needed two signatures for each key; ShortSig needed five and six.  All four
fresh-message forgeries were accepted at the exact specification sizes of 3,124
and 1,620 bytes after the original secret and collected transcript state were erased.
Changed-message, response-tamper, and digest-tamper controls all rejected.

For every extracted edge, both implementations compare the complete canonical
tensors and verify the corresponding public tensor-equivalence equation, not
merely a canonical digest.  The Balanced run then checks its single recovered
base-to-nonbase map.  The ShortSig run checks every edge and every composed
base-to-label map before using connected-graph completion.  These runs
establish collision observability, certificate orientation, graph composition,
witness sufficiency, verifier binding, and ordinary fresh randomness.  They do
not estimate or execute the full 128-bit-domain attack; those costs come from
the separately qualified analytic models.

\section{Specification and implementation boundary}

The untouched submitted source expands each round from
\begin{equation}
 \texttt{"MEDS2END-TRINE-ROUND-v1"}\mathbin\Vert
 \mathit{salt}\mathbin\Vert s_i\mathbin\Vert
 \mathsf{LE32}(i),                                      \label{eq:sourcedomain}
\end{equation}
where $\mathit{salt}$ is a fresh serialized $2\lambda$-bit value.  Its level-I
signatures are therefore 32 bytes longer than the specified values in
Table~\ref{tab:params}.  Independent salts make the decoder domain
signature-specific, so equality of the serialized round seed no longer gives
equality of the expanded point.  The ordinary attack in this report does not
apply to that source construction: the untouched submitted C source is not
vulnerable to this accumulation mechanism.  This report is not an
implementation-break claim.  It targets the direct unsalted,
message-independent realization described in the specification, while the
source already contains the relevant salt-based repair.

The reproducer contains two patches.  The minimal conformance patch removes
only the source-only salt from sampling, serialization, codecs, and the round
XOF while preserving the domain string, round index, SHAKE256, field sampler,
$\Corank$ retry logic, canonicalizer, challenge parser, and verifier.  A second
patch adds only reduced-seed and certificate instrumentation for the hostile
experiment.

The specification leaves the precise seed decoder implicit.  A contrary interpretation
could define $D_i(s,M)$, binding the message without adding signature bytes.
That construction would defeat cross-message accumulation.  It is not the
written order of Algorithm~6, where $a_i\sample\Corank(\phi_X)$ precedes the
message hash and has no message argument, and it is not the submitted decoder.
For this reason, the exact claim is an ordinary collision attack on the
\emph{direct unsalted, message-independent realization described in the
specification}.  If the submitted
salt is declared normative, or a message-bound decoder is specified, the
claim does not apply.  Nothing in the analysis treats a source deviation as a
vulnerability; the minimal patch is used only to test the construction
described in the specification that the paper analyzes.

\section{Repair}

The specification should require a per-signature public salt of at least $2\lambda$ bits,
serialize it, include it in the signature-size formula, and bind it together
with the round index and full instance identifier before any round expansion.
For example,
\begin{equation}
 D_{\mathit{inst},i,\rho}(s)=
 \mathsf{Decode}\bigl(
   \mathsf{XOF}(\mathsf{domain}\mathbin\Vert\mathit{inst}
   \mathbin\Vert\rho\mathbin\Vert s\mathbin\Vert\mathsf{LE32}(i))
 \bigr).                                                \label{eq:repair}
\end{equation}
The source already implements the central salt-and-round part of this repair.
Alternatively, binding the complete message digest into the decoder prevents
cross-message accumulation, but the decoder and proof must then state that
dependency explicitly.  Tests should assert that equal serialized round seeds
under two independently salted signatures produce unrelated commitments.

\section{Prior art and novelty boundary}

Ward Beullens's public note \emph{Trivial multi-key attacks + attack on ALTEQ},
posted to the NIST PQC Forum on 21 February 2024, describes the generic public-
table collision/guessing attack on $\lambda$-bit signing seeds and recommends
deriving rounds from the seed together with a salt and round identifier
\cite{beullens2024}.  Arnaud Sipasseuth's \emph{DVA: Dangerous Variations of
ALTEQ}, IACR ePrint 2024/817, Section~3.2, explicitly includes internal seed
collisions and the $\mathit{salt}\mathbin\Vert\mathit{rid}$ repair
\cite{sipasseuth2024}.  They own the generic mechanism and repair.

The contribution claimed here is narrower: applying that mechanism to the
unsalted construction described in the TRINE specification; identifying the
discrepancy with the already salted source;
using TRINE's public canonicalization certificate to extract a full tensor
equivalence; showing that Balanced needs one edge while ShortSig needs a
connected five-label graph; deriving the exact level-I tradeoffs with mandatory
$\Corank$ retry costs; and validating fresh forgeries at full TRINE dimensions
with reduced seed entropy.  The connected-graph repetition is a
candidate-specific completion of the known per-label attack, not a new generic
attack family.

A best-effort check of the public NGCC TRINE page on 3 October 2026 found only
report \texttt{sign-30-1}, an unrelated deterministic-key implementation
failure~\cite{ngcc-trine}.  No public TRINE-specific report of this attack on
the unsalted specification's seed domain was found at that checkpoint.  This
is a dated best-effort negative search, not proof of novelty or absence.

\section{Reproduction}

The companion ZIP is self-contained and contains the specification-aligned source,
minimal conformance patch, hostile instrumentation patch, Balanced and ShortSig
attack harnesses, portable controls, resource model, frozen outputs, validator,
and checksum manifest.  From its \texttt{reproducer} directory, run
\begin{center}
                 \texttt{./run\_all.sh}.
\end{center}
The two full-dimension reduced-seed profile runs execute concurrently.  On the
audit host, the slower frozen run took approximately eight minutes and used
about 72 MiB peak resident memory.  The command ends with
\texttt{validation=PASS}.  To check the sealed evidence and manifest without
rerunning the native experiments, use
\texttt{./run\_all.sh --validate-only}.  No command depends on an audit-cache or
private workspace path.

\section{Conclusion}

The direct unsalted construction described in the TRINE specification places
every base commitment in a $2^{128}$ seed domain.  Ordinary independent
signatures eventually collide at
one round index.  Different challenge labels turn that birthday event into the
canonical-form equality that TRINE's own public certificates extract.  One
edge completes Balanced; a connected label graph completes ShortSig.  The
resulting equivalent witness creates ordinary fresh-message signatures.

The attack algorithms cover the complete extraction-to-forgery chain and give
full-parameter cost estimates, but a full-entropy forgery was not physically
executed.  In particular, the exact fixed-index point at $Q=2^{64}-1$ has
approximately one-half ideal-model success and is not excluded by any
specified TRINE per-key lifetime cap.  Its measured attacker work is about
$2^{94.25}$ and $2^{96.99}$ cycles, and signer work, traffic, and storage are
reported separately.  Other Pareto points reduce attacker work or query data
at different physical costs.  Only reduced-seed end-to-end forgeries were
run.  The untouched source is protected by a fresh serialized
$2\lambda$-bit salt, is not vulnerable to this accumulation mechanism, and is
outside the claim.  The generic attack belongs to Beullens and Sipasseuth; the
result here is the concrete TRINE instantiation of the unsalted specification
and its validated extraction-to-forgery chain.

\section*{Acknowledgements}

The computations were performed on the Norwegian Research and Education Cloud
(NREC), using resources provided by the University of Bergen and the University
of Oslo.

\begin{thebibliography}{9}
\small

\bibitem{trine-spec}
Gang Tang, Debiao He, Yu Dai, Cong Peng, Min Luo, Yinan Li, Xinyi Huang,
Runqing Xu, Zihao Yu, and Bingbing Xia,
\emph{The TRINE Signature Scheme: Algorithm Specifications and Supporting
Documentation}, NGCC submission, 2026.

\bibitem{beullens2024}
Ward Beullens,
\emph{Trivial multi-key attacks + attack on ALTEQ},
NIST PQC Forum, 21 February 2024,
\url{https://groups.google.com/a/list.nist.gov/g/pqc-forum/c/tjhrrmv837w/m/sjxHooYgBAAJ}.

\bibitem{sipasseuth2024}
Arnaud Sipasseuth,
\emph{DVA: Dangerous Variations of ALTEQ},
IACR Cryptology ePrint Archive, Report 2024/817, 2024,
\url{https://eprint.iacr.org/2024/817}.

\bibitem{ngcc-trine}
NGCC Security Reports,
\emph{TRINE (sign-30)},
\url{https://ngcc.dev/reports/sign-30.html}, accessed 3 October 2026.

\end{thebibliography}

\end{document}
